% !TEX program = pdflatex
% compilar con pdflatex (2 pasadas) desde esta carpeta (informe/)
\documentclass[11pt,a4paper]{article}

\usepackage[utf8]{inputenc}
\usepackage[T1]{fontenc}
\usepackage[spanish,es-tabla]{babel}
\usepackage[margin=2.3cm]{geometry}
\usepackage{amsmath,amssymb}
\usepackage{booktabs}
\usepackage{enumitem}
\usepackage{xcolor}
\usepackage{float}
\usepackage[colorlinks=true,linkcolor=blue!60!black,citecolor=blue!60!black,urlcolor=blue!60!black]{hyperref}

\newcommand{\B}{\mathbf{B}}
\newcommand{\vct}{\mathbf{v}}
\newcommand{\dif}{\partial_t}
\newcommand{\code}[1]{\texttt{#1}}
\newcommand{\term}[1]{\textbf{#1}}

\setlength{\parskip}{2pt}
\setlength{\emergencystretch}{2em}

\title{\textbf{Proyecto 2 — Glosario razonado: conceptos, términos y ecuaciones}\\[4pt]
\large Compañero de estudio de la guía 2D (\code{Rayleigh\_Taylor\_MHD/}), del informe 3D
(\code{SNR\_3D/informe/}) y del notebook de análisis}
\author{Zeuz Espinosa\\ Astrof\'isica Computacional 2026-II}
\date{Noviembre de 2026}

\begin{document}
\maketitle

\begin{quote}\small
\textbf{Para qué es esto.} Documento de consulta para entender \emph{a cabalidad} lo que se
hizo en el proyecto: cada concepto físico, cada término numérico y cada ecuación que aparece
en la guía 2D, en la corrida PLUTO 3D del SNR, en el notebook o en el informe. Está organizado
de lo físico a lo numérico: la Parte~I es la mecánica de fluidos magnetizada, la Parte~II la
inestabilidad, la Parte~III el contexto de remanentes de supernova, la Parte~IV los métodos
numéricos y los diagnósticos, y la Parte~V es un glosario alfabético rápido para cuando solo
se necesita recordar qué significa una sigla. Los apéndices traen la tabla de símbolos y el
mapa ``ecuación $\to$ dónde se usa en el proyecto''.
\end{quote}

\tableofcontents
\clearpage

%=====================================================================
\part*{Parte I --- El fluido y el campo: mecánica de fluidos magnetizada}
\addcontentsline{toc}{section}{Parte I --- El fluido y el campo}

\section{Qué es la MHD y qué asume}

\term{Magnetohidrodinámica (MHD)}: tratar un plasma (gas ionizado) como un \emph{fluido}
conductor que arrastra y es empujado por campos magnéticos. La versión \term{MHD ideal}
que usa PLUTO en este proyecto asume tres cosas:

\begin{enumerate}[nosep]
  \item \textbf{Casi-neutralidad y un solo fluido}: electrones e iones están tan acoplados
  por colisiones/campos que se mueven juntos con una sola velocidad $\vct$. Las escalas
  de interés (la caja, los dedos) son enormes comparadas con el radio de Larmor y el
  camino libre medio.
  \item \textbf{Conductividad infinita} ($\eta_\mathrm{res}=0$, sin resistividad): de la ley
  de Ohm ideal $\mathbf{E}+\vct\times\B=0$ sigue que el campo está \term{congelado} al
  fluido (teorema de Alfvén): el flujo magnético a través de cualquier contorno que se
  mueve con el plasma se conserva. Esto explica por qué el estiramiento del flujo
  \emph{amplifica} el campo en las puntas de los dedos (¡es lo que mide una de las figuras!).
  \item \textbf{Corriente de desplazamiento despreciable} (Ampère no relativista):
  $\mathbf{J}=\nabla\times\B$ (en unidades de código; en Gaussianas
  $\mathbf{J}=c\,\nabla\times\B/4\pi$).
\end{enumerate}

Cuándo es válida: escalas $L \gg$ camino libre medio y $L \gg$ radio de Larmor, y frecuencias
$\ll$ las de plasma. En un SNR joven, en el medio interestelar y en prominencias solares se
cumple de sobra; en la magnetosfera o a escalas menores que la inercia iónica habría que
usar Hall o dos fluidos.

\section{Las cinco ecuaciones del proyecto, en palabras}

El sistema que resuelve PLUTO (unidades de código, ver \S\ref{sec:unidades}) es:

\begin{align}
\underbrace{\dif{\rho} + \nabla\cdot(\rho\vct)}_{\text{masa}}
&= 0 \\[3pt]
\underbrace{\dif{(\rho\vct)} + \nabla\cdot\Big[\rho\vct\vct + \big(p+\tfrac{B^2}{2}\big)\mathbf{I} - \B\B\Big]}_{\text{momento}}
&= \underbrace{\rho\,\mathbf{g}}_{\text{gravedad}} \\[3pt]
\underbrace{\dif{E} + \nabla\cdot\Big[\big(E+p+\tfrac{B^2}{2}\big)\vct - (\B\cdot\vct)\B\Big]}_{\text{energía total}}
&= \underbrace{\rho\,\vct\cdot\mathbf{g}}_{\text{trabajo gravitacional}} \\[3pt]
\underbrace{\dif{\B} - \nabla\times(\vct\times\B)}_{\text{inducción}}
&= 0,
\qquad \underbrace{\nabla\cdot\B = 0}_{\text{constraint}}
\end{align}
con $\dif\equiv\partial/\partial t$, la ecuación de estado ideal
\begin{equation}
p = (\gamma-1)\,e_\mathrm{int},
\qquad\text{equivalentemente}\qquad
E \;=\; \underbrace{\frac{p}{\gamma-1}}_{\text{interna}}
+\underbrace{\frac{\rho v^2}{2}}_{\text{cinética}}
+\underbrace{\frac{B^2}{2}}_{\text{magnética}},
\end{equation}
con $\gamma=5/3$ (gas monoatómico, grados de libertad traslacionales).

Qué dice cada una, en palabras:
\begin{itemize}[nosep]
  \item \textbf{Masa}: la densidad de un elemento de fluido solo cambia si el flujo lo
  comprime/expande. Sin fuentes: la masa total se conserva exactamente.
  \item \textbf{Momento}: es $\mathbf{F}=m\mathbf{a}$ escrito para un volumen de fluido. El
  flujo transporta momento por advección ($\rho\vct\vct$), por presión de gas ($p\mathbf{I}$)
  y por el campo ($+$ y $-$ términos magnéticos, ver \S\ref{sec:maxwell}). El lado derecho
  es la fuente gravitacional: en nuestro problema $\mathbf{g}=g\,\hat{\mathbf{y}}$ con
  $g<0$ constante (el \code{GRAV} del \code{pluto.ini}).
  \item \textbf{Energía}: la energía \emph{total} por volumen $E$ se conserva salvo por el
  trabajo que hace la gravedad, $S_E=\rho\vct\cdot\mathbf{g}$. Con este término incluido,
  en un dominio cerrado (fronteras reflectivas) la suma
  $E_\mathrm{int}+E_\mathrm{cin}+E_\mathrm{mag}+E_\mathrm{pot}$ es constante: por eso en el
  notebook la deriva medida dio $\sim10^{-11}$.
  \item \textbf{Inducción}: es la ley de Faraday con la Ohm ideal sustituida; dice cómo el
  movimiento curva y estira el campo. Junta con el congelamiento, es la ecuación que
  ``fabrica'' el $B_y$ que no existía al inicio.
  \item \textbf{Constraint}: $\nabla\cdot\B=0$ no es una ecuación a resolver sino una
  \emph{condición de compatibilidad}: si se cumple en $t=0$, las otras ecuaciones la
  preservan en la solución exacta. Numéricamente NO es automática (\S\ref{sec:divb}).
\end{itemize}

\section{Presión magnética y tensión magnética: la clave de todo}
\label{sec:maxwell}

La fuerza por unidad de volumen que ejerce el campo sobre el fluido es (unidades de código;
en Gaussianas dividir por $4\pi$):
\begin{equation}
\mathbf{f}_\mathrm{mag} = \mathbf{J}\times\B = -\nabla\Big(\frac{B^2}{2}\Big)
+ (\B\cdot\nabla)\B
\qquad (\text{si } \nabla\cdot\B=0).
\end{equation}
Los dos términos tienen interpretaciones físicas distintas que se usan todo el tiempo en
este proyecto:

\begin{itemize}
  \item \textbf{Presión magnética}, $-\nabla(B^2/2)$: el campo empuja \emph{perpendicular}
  a sus líneas igual que un gas de presión $B^2/2$. Por eso la presión total que aparece en
  el tensor de momento es $p+B^2/2$. Cuando el campo se amplifica por estiramiento, esta
  presión crece y termina frenando la compresión de las puntas.
  \item \textbf{Tensión magnética}, $(\B\cdot\nabla)\B$: las líneas de campo se comportan
  como cuerdas de guitarra tensas; si se curvan, tiran de regreso para enderezarse. Es
  \emph{exactamente} la fuerza que estabiliza los modos cortos de la RT magnetizada
  (\S\ref{sec:rtmhd}): una línea que atraviesa la interfaz corrugada queda estirada y
  ``se resiste'' a que la corrugación sea más fina que su longitud natural de curvatura.
\end{itemize}

Dos palabras que aparecen en la literatura del proyecto y vienen de aquí:
\term{drapeo} (el campo ambiental se ``envuelve'' alrededor de la superficie de choque /
de los dedos, concentrándose en sus flancos) y \term{congelamiento} (el fluido arrastra las
líneas; estirar el fluido estira y amplifica el campo: $\frac{D}{Dt}\Big(\frac{\B}{\rho}\Big)=\Big(\frac{\B}{\rho}\cdot\nabla\Big)\vct$ para un elemento de fluido).

\section{Velocidades características, ondas y el parámetro $\beta$}

En un fluido compresible las perturbaciones viajan como ondas. La \term{velocidad del
sonido} sale de ``cuánto sube la presión al comprimir'': para un gas ideal adiabático,
\begin{equation}
c_s^2=\Big(\frac{\partial p}{\partial\rho}\Big)_s=\frac{\gamma p}{\rho}.
\end{equation}
En el setup, $p_0=1/\gamma$ y $\rho_1=1$ se escogen justamente para que $c_s=1$ en el fluido
ligero: la velocidad del sonido es la unidad de velocidad del problema.

Con campo, hay tres familias de ondas:
\begin{itemize}[nosep]
  \item \term{Ondas de Alfvén}: transversales, viajan \emph{a lo largo} del campo a la
  \term{velocidad de Alfvén} $v_A=B/\sqrt{\rho}$ (en Gaussianas $B/\sqrt{4\pi\rho}$).
  Físicamente: la tensión magnética es la fuerza restauradora, como la tensión de una
  cuerda de densidad $\rho$.
  \item \term{Ondas magnetosónicas rápidas y lentas}: compresionales, una mezcla de
  presión de gas y presión magnética. La \term{rápida} es la que importa para el paso de
  tiempo; su velocidad (en la dirección $\hat{\mathbf{n}}$) es
  \begin{equation}
  c_f^2=\frac{c_s^2+v_A^2}{2}
  +\sqrt{\Big(\frac{c_s^2+v_A^2}{2}\Big)^2 - c_s^2\,v_{A,n}^2}\,,
  \qquad v_{A,n}^2=\frac{B_n^2}{\rho}\,,
  \end{equation}
  y satisface $\max(c_s, v_{A,n})\le c_f\le\sqrt{c_s^2+v_A^2}$. Casos límite útiles:
  propagación a lo largo del campo: $c_f=\max(c_s,v_A)$; perpendicular: $c_f=\sqrt{c_s^2+v_A^2}$.
\end{itemize}

La \term{relación de presiones plasmática} (o \emph{beta plasmática}) compara gas contra
campo:
\begin{equation}
\beta \equiv \frac{p_\mathrm{gas}}{p_\mathrm{mag}} = \frac{2p}{B^2}
\qquad\text{(Gaussianas: } 8\pi p/B^2\text{)}.
\end{equation}
$\beta\gg1$ (nuestro caso: $\sim10^3$--$10^4$): el campo es débil en presión; $\beta\ll1$:
domina el campo. \emph{Trampa conceptual del proyecto}: aunque $\beta\gg1$, el campo controla
la RT porque actúa por \emph{tensión} sobre las escalas cortas de la interfaz, no por
presión global. Débil $\neq$ irrelevante.

\section{El problema de $\nabla\cdot\B=0$ y sus dos soluciones}
\label{sec:divb}

Si el esquema numérico deja un residuo $\nabla\cdot\B\neq0$, ese ``monopolo'' espurio
produce fuerzas no físicas (el término $\mathbf{J}\times\B$ con $\mathbf{J}=\nabla\times\B$
deja de ser perpendicular a $\B$), transporte erróneo de momento y hasta inestabilidades.
Las dos estrategias usadas en el proyecto:

\begin{itemize}
  \item \term{Transporte restringido} (CT, \emph{constrained transport}): la malla es
  \emph{escalonada} (\term{staggered}): $B_x$ vive en las caras $x$ de las celdas,
  $B_y$ en las caras $y$, etc. La inducción se escribe en forma integral (Faraday):
  \begin{equation}
  \frac{d}{dt}\int_\mathrm{cara}\B\cdot d\mathbf{S} = -\oint_\mathrm{borde}\mathbf{E}\cdot d\mathbf{l}
  \end{equation}
  con el campo eléctrico inducido $\mathbf{E}=-\vct\times\B$.
  y el flujo por una cara se actualiza con la \term{FEM} (fuerza electromotriz, en 2D la
  componente $E_z$) definida en sus bordes/esquinas. La divergencia discretizada es la suma
  de los cuatro flujos de las caras alrededor de cada celda: como cada FEM entra con signos
  opuestos en las dos celdas vecinas, \emph{se cancela idénticamente}: $\nabla\cdot\B=0$ a
  precisión de máquina ($\sim10^{-15}$). Es lo que usa la config \#07 (2D) que recomienda la
  guía.
  \item \term{Limpieza GLM (\emph{divergence cleaning} de Dedner)}: se agrega una variable
  escalar $\psi$ y un par de ecuaciones acopladas que convierten el error de divergencia en
  una onda que se propaga rápido (a $c_h\sim$ velocidad máxima del sistema) y se amortigua
  (tiempo $\tau_d$): aproximadamente
  \begin{equation}
  \dif{\B}+\nabla\psi=\nabla\times(\vct\times\B),\qquad
  \dif{\psi}+c_h^2\nabla\cdot\B=-\psi/\tau_d .
  \end{equation}
  El error no se elimina: se ``barrre y disipa''. Es lo que usa la config \#09 3D de la
  corrida del SNR; por eso el notebook midió un residuo relativo $\sim0.2$ (en unidades de
  $B_0/\Delta x$, contra el campo inicial) y no $10^{-15}$. Ese es el precio normal de GLM.
\end{itemize}

%=====================================================================
\part*{Parte II --- La inestabilidad}
\addcontentsline{toc}{section}{Parte II --- La inestabilidad}

\section{Rayleigh--Taylor hidrodinámica: por qué crece y con qué tasa}

\term{Configuración}: fluido pesado ($\rho_2$) \emph{arriba}, ligero ($\rho_1$) abajo,
aceleración efectiva de magnitud $g$ apuntando ``hacia abajo'' ($-\hat{\mathbf{y}}$).
Es inestable: es el vaso con agua boca abajo, el hongo de la bomba, los dedos del SNR.
La condición general: la aceleración apunta del fluido ligero al pesado.

\subsection{El número de Atwood}

\begin{equation}
A \equiv \frac{\rho_2-\rho_1}{\rho_2+\rho_1}\in(0,1]\,.
\end{equation}
Mide la asimetría de densidades. $A\to0$: casi simétrico (burbujas y picos iguales);
$A\to1$: burbuja en vacío (el pico cae casi en caída libre). En la corrida 2D de la guía
$A=1/3$ ($\eta=2$); en la 3D del SNR $A=3/5$ ($\eta=4$).

\subsection{El mecanismo: vorticidad baroclínica}

En el equilibrio, $\nabla p$ es vertical y $\nabla\rho$ es normal a la interfaz. Mientras
la interfaz es plana, ambos son paralelos. Al corrugarse, dejan de serlo: en los flancos de
la corrugación aparece el término \term{baroclínico} de la ecuación de la vorticidad
($\boldsymbol{\omega}=\nabla\times\vct$):
\begin{equation}
\frac{D\boldsymbol{\omega}}{Dt}
= (\boldsymbol{\omega}\cdot\nabla)\vct
-\boldsymbol{\omega}\,\nabla\cdot\vct
+ \underbrace{\frac{\nabla\rho\times\nabla p}{\rho^2}}_{\text{baroclínico}}\,.
\end{equation}
Esa vorticidad generada es del signo que hace subir la cresta y bajar el valle: la
corrugación se amplifica sola. La energía sale del campo gravitacional: al mezclarse, el
centro de masa baja y la energía potencial se libera como movimiento.

\subsection{La tasa de crecimiento, deducida}

Asumimos incompresible ($\nabla\cdot\vct=0$), sin viscosidad, perturbación
$\propto e^{ikx+\sigma t}$ de la interfaz $\xi(x,t)$, fluidos semi-infinitos. Como en el
volumen no hay vorticidad inicialmente ni baroclínica fuera de la interfaz, el flujo en el
interior de cada fluido es potencial: $\vct=\nabla\phi$, y la incompresibilidad da
$\nabla^2\phi=0$, con soluciones que decaen alejándose de la interfaz:
\begin{equation}
\phi_1=\Phi_1\,e^{ikx+ky+\sigma t}\ (y<0),
\qquad
\phi_2=\Phi_2\,e^{ikx-ky+\sigma t}\ (y>0).
\end{equation}
Ecuación de Bernoulli linealizada en cada fluido:
$\rho_i\,\sigma\phi_i + \delta p_i + \rho_i\,g\,y = 0$. Dos condiciones en la interfaz
desplazada $y=\xi$:
\begin{enumerate}[nosep]
  \item \textbf{Cinemática} (el fluido sigue a la interfaz):
  $\dif{\xi}=v_y=\partial_y\phi|_{0}$ $\Rightarrow$ $\sigma\xi=k\Phi_1=-k\Phi_2$.
  \item \textbf{Continuidad de presión} en $y=\xi$: la presión perturbada de cada lado más
  el desplazamiento del gradiente hidrostático deben coincidir:
  $\delta p_1-\delta p_2 = -(\rho_2-\rho_1)\,g\,\xi$.
\end{enumerate}
Sustituyendo $\delta p_i=-\rho_i\sigma\phi_i$ y usando la cinemática para eliminar
$\Phi_i$ a favor de $\xi$ queda, tras dos líneas de álgebra,
\begin{equation}
\boxed{\ \sigma^2 = A\,g\,k\ } \qquad
(\text{crecimiento exponencial } e^{\sigma t},\ \sigma=\sqrt{Agk}).
\label{eq:sigma-hd}
\end{equation}
Lecturas útiles de la fórmula: modos más cortos (mayor $k$) crecen más rápido
($\sigma\propto\sqrt{k}$, sin corte en esta teoría); mayor contraste (mayor $A$) crece más
rápido; en hidrodinámica pura no hay escala de longitud preferida --- el corte lo pone la
viscosidad o, en la simulación, la difusión numérica (y en la versión magnetizada, la
tensión del campo).

Con una semilla en velocidad de amplitud $v_0$ (como la del setup), la solución exacta de
la fase lineal es
\begin{equation}
\xi(t)=\xi_0\cosh(\sigma t)+\frac{v_0}{\sigma}\sinh(\sigma t)
\;\simeq\;\frac{v_0}{2\sigma}e^{\sigma t}\quad (\sigma t\gtrsim2),
\end{equation}
de donde sale la ``amplitud inicial efectiva'' $\xi_{0,\mathrm{eff}}=v_0/2\sigma$ que usa la
guía para estimar cuándo arranca el régimen no lineal.

\section{Fases no lineales: burbujas, picos y Layzer}

Cuando $\xi\sim\lambda$, el crecimiento exponencial muere: la burbuja asciende a velocidad
\emph{constante} y el pico cae cada vez más rápido (a velocidad casi constante si $A$
moderado, tendiendo a caída libre si $A\to1$). El modelo clásico (\term{Layzer 1955};
generalizado a $A$ arbitrario por \term{Goncharov 2002}) aproxima la burbuja por flujo
potencial dentro de un tubo de ancho $\lambda$ y da, en 2D,
\begin{equation}
v_b=\sqrt{\frac{2A\,g\,\lambda}{\pi(1+A)}}\,.
\end{equation}
En los flancos de burbujas y picos hay cizalla $\rightarrow$ se disparan inestabilidades
\term{Kelvin--Helmholtz (KH)} secundarias, que son los ``rulos'' de los hongos clásicos.
Una perturbación \term{multimodo} (como la aleatoria de la corrida 3D) produce una capa de
mezcla cuyo espesor crece primero exponencialmente y luego $\propto t$ (burbujas compitiendo
por el ascenso, fusión de modos): eso es lo que mide $h_{95-5}(t)$ en el notebook.

\section{RT magnetizada: la ecuación de Chandrasekhar, deducida}
\label{sec:rtmhd}

Ahora $\B=B\,\hat{\mathbf{x}}$ uniforme a ambos lados. Tres ideas nuevas:
\begin{enumerate}[nosep]
  \item La inducción linealizada con $\nabla\cdot\vct=0$ y $\B$ uniforme da
  $\sigma\,\delta\B=(\B\cdot\nabla)\vct \Rightarrow \delta\B=ikB\vct/\sigma$: la
  perturbación del campo es proporcional a la de velocidad.
  \item En el momento linealizado aparece la tensión:
  $\rho\sigma\vct=-\nabla\delta P+(\B\cdot\nabla)\delta\B$ con $\delta P$ la presión
  \emph{total} (gas $+$ magnética). Combinando con 1), la tensión actúa como una
  ``densidad extra'': $\delta P=-(\rho\sigma+k^2B^2/\sigma)\phi$.
  \item Las condiciones de frontera en la interfaz son las mismas que en \S previo
  (cinemática y continuidad de presión total, ahora incluyendo la perturbación de
  presión magnética que ya quedó absorbida en el término $k^2B^2/\sigma$).
\end{enumerate}
Repitiendo el álgebra del caso hidrodinámico (la cuenta completa está en la guía 2D,
\S5.2) llega la \term{relación de dispersión de Chandrasekhar} (unidades de código;
en Gaussianas el término magnético lleva un $1/4\pi$):
\begin{equation}
\boxed{\ \sigma^2 = A\,g\,k \;-\; \frac{2\,(\mathbf{k}\cdot\B)^2}{\rho_1+\rho_2}\ }
\label{eq:sigma-mhd}
\end{equation}
El factor 2 no es decorativo: la tensión restauradora la ejercen las líneas de \emph{ambos}
lados de la interfaz (cada fluido aporta su $k^2B^2$). Intuición de cuerda: desplazar la
interfaz transversalmente estira las líneas como cuerdas de densidad $\rho_1$ y $\rho_2$;
el modo corta cuando la tensión gana a la flotabilidad.

De la eq.~\eqref{eq:sigma-mhd} salen \emph{todas} las cantidades que aparecen en el
proyecto:

\begin{itemize}
  \item \textbf{Anisotropía}: solo el componente $\mathbf{k}\parallel\B$ ($k_x$) entra en
  el término magnético. Modos \term{de intercambio} ($\mathbf{k}\perp\B$, o sea $k_z$)
  crecen como en hidrodinámica pura; modos \term{undulares} ($\mathbf{k}\parallel\B$) sienten
  la tensión. En 2D (guía) todo modo es undular: supresión máxima. En 3D (SNR) hay toda la
  gama de ángulos: la mezcla queda elongada a lo largo del campo. Es la base del diagnóstico
  de sectores del notebook.
  \item \textbf{Longitud de corte}: con $\mathbf{k}\parallel\B$,
  $\sigma^2>0$ solo si $\lambda>\lambda_c$. Con $B$ escrito en unidades del campo crítico
  ($B^2=\chi^2(\rho_2-\rho_1)gL_x/4\pi$, ver abajo),
  \begin{equation}
  \sigma^2(\lambda) = A g \frac{2\pi}{\lambda}\Big(1-\chi^2\frac{L_x}{\lambda}\Big)
  \quad\Rightarrow\quad
  \lambda_c = \chi^2 L_x\,,
  \end{equation}
  modos paralelos al campo con $\lambda<\lambda_c$ quedan estabilizados.
  \item \textbf{Campo crítico}: $B_c\equiv\sqrt{(\rho_2-\rho_1)gL_x}$ (la definición que
  usa el \code{init.c}; en unidades de código el campo inicial es $\chi B_c/\sqrt{4\pi}$).
  $\chi=1$ estabiliza \emph{exactamente} el modo más largo que cabe en el dominio
  ($\lambda=L_x$); $\chi>1$ congela todo. En la guía 2D esto da la predicción estrella
  $\sigma_\mathrm{fund}=\sqrt{Agk}\sqrt{1-\chi^2}$.
  \item \textbf{Modo de máximo crecimiento}: derivando respecto de $k$ (a lo largo de
  $\B$), el máximo está en $k_m=k_c/2$, es decir $\lambda_m=2\lambda_c$, con
  $\sigma_\mathrm{max}=\sqrt{A g\pi/(2\chi^2L_x)}$.
  \item \textbf{Con $\mathbf{k}\perp\B$ no hay corte} en esta teoría incompresible: $\sigma$
  crece como $\sqrt{k}$ sin freno magnético, y el límite real lo pone la malla (Nyquist,
  \S\ref{sec:fourier}). En la corrida 3D: $\lambda_c=0.04$ son solo $\sim2$ celdas --- el
  corte magnético duro está apenas fuera de lo que la malla resuelve, y la anisotropía
  que se ve es la de \emph{tasas} (P$_\perp$/P$_\parallel\approx17$ medida).
\end{itemize}

%=====================================================================
\part*{Parte III --- El contexto astrofísico: remanentes de supernova}
\addcontentsline{toc}{section}{Parte III --- El contexto astrofísico}

\section{La vida de un remanente}

Un \term{remanente de supernova (SNR)} es lo que queda de la onda de choque de la
explosión barriendo el medio circundante, junto con los restos de la estrella
(\term{eyecta}). Su evolución tiene fases:
\begin{enumerate}[nosep]
  \item \term{Expansión libre}: la eyecta es mucho más densa que el medio; $R\propto t$
  ($\eta=1$) hasta barrear una masa comparable a la propia.
  \item \term{Fase de eyecta dominada / Sedov--Taylor}: el choque delantero es fuerte y
  la energía se conserva; la solución de \term{Sedov} da
  $R\propto (Et^2/\rho_0)^{1/5}$, o sea $\eta=2/5$ (más la cola de la eyecta con perfil
  en ley de potencias detrás: descripción auto-semejante de \term{Chevalier 1982}). Es la
  fase de Tycho y Cas~A, y la de nuestra simulación.
  \item \term{Fase radiativa}: el gas enfría por radiación, el choque pierde energía,
  $R\propto t^{2/7}$ aprox., el SNR se dispersa en el medio interestelar.
\end{enumerate}
Dentro del choque delantero hay dos superficies importantes: el propio \term{choque}
(salto de presión, velocidad y densidad; comprime y calienta) y la \term{discontinuidad
de contacto} entre eyecta chocada y medio chocado (misma presión a ambos lados,
\emph{densidades distintas}: solo hay que empujar densidad a través, no momento --- por eso
es precisamente donde vive la RT).

\section{Por qué el contacto es Rayleigh--Taylor: la gravedad efectiva}

Si $R(t)=C\,t^\eta$ con $\eta<1$, la superficie desacelera:
\begin{equation}
a = \ddot{R} = \eta(\eta-1)\,\frac{R}{t^2}
\;\;\Rightarrow\;\;
|a| = \eta(1-\eta)\,\frac{R}{t^2}\,.
\end{equation}
(Con Sedov, $\eta=2/5$: $|a|=\tfrac{6}{25}R/t^2$.) Orden de magnitud para Tycho
($R\sim5$ pc, $t\sim430$ yr): $|a|\sim2\times10^{-4}\ \mathrm{m/s^2}$ --- ridícula como
gravedad terrestre, pero suficiente: en el marco que desacelera con la interfaz, hay una
aceleración ``hacia adentro'' que apunta del medio chocado (ligero) a la eyecta chocada
(pesada). Condición de Rayleigh--Taylor: cumplida. Resultado: el contacto se rompe en
dedos de eyecta que penetran hacia afuera --- exactamente los picos que caen en la
simulación, con la correspondencia
\begin{center}
\emph{fluido pesado arriba = eyecta chocada;\ \ fluido ligero abajo = medio chocado;\
$g$ = desaceleración;}
\\ \emph{$\B$ horizontal = campo ambiental drapado en el choque delantero.}
\end{center}
En el SNR real la aceleración decae con el tiempo y la geometría es esférica; el parche
local plano con $g$ constante es la aproximación de trabajo estándar (la misma de la
literatura del proyecto: Stone \& Gardiner 2007 para el método, Orlando et al.\ 2012 y
Fraschetti et al.\ 2010 para el caso SNR).

%=====================================================================
\part*{Parte IV --- El lado numérico: todo lo que aparece en PLUTO y en el notebook}
\addcontentsline{toc}{section}{Parte IV --- El lado numérico}

\section{Unidades de código: qué significa ``$c_s=1$''}
\label{sec:unidades}

PLUTO no simula en unidades físicas sino en \term{unidades de código}: cantidades
adimensionales donde uno fija qué vale ``1''. El setup normaliza
\emph{longitud} con el ancho horizontal $L_x=1$, \emph{densidad} con el fluido ligero
$\rho_1=1$, y \emph{presión} con $p_0=1/\gamma$ de modo que la velocidad del sonido del
fluido ligero sea $c_s=1$; de ahí salen las unidades derivadas de tiempo ($L/c_s$),
velocidad y campo ($B^2\to\rho v^2$). Traducir a física real es escoger tres números de
escala ($L$, $\rho$, $c_s$) y multiplicar. El factor $1/\sqrt{4\pi}$ en el \code{init.c}
al inicializar $B$ viene de pasar de la convención CGS del campo crítico a la convención
sin $4\pi$ de las ecuaciones del código.

\section{Volúmenes finitos y el método de Godunov}

El dominio se parte en \term{celdas}; en cada una se guardan \term{promedios de volumen}
de las variables \term{conservativas} $U=(\rho,\rho\vct,E,\B)$. Evolucionar es calcular
los \term{flujos} por las caras: $\dif{U} + \nabla\cdot\mathbf{F}(U)=S$. Ventaja: lo que
sale de una celda entra en la vecina --- conservación exacta, clave para choques.

¿Flujo por una cara donde el estado izquierdo y derecho difieren? Eso es un
\term{problema de Riemann}, y el \term{solver de Riemann} lo resuelve (aproximadamente):
\begin{itemize}[nosep]
  \item \term{Roe}: linealiza el problema (matriz jacobiana con un estado promedio de
  Roe) y resuelve exacto el problema linealizado. Es el solver de la corrida 3D
  (\code{Solver = roe}).
  \item \term{HLL}: solo usa las velocidades de las ondas más rápida y más lenta
  ($S_L$, $S_R$) y asume un estado único en el abanico:
  \begin{equation}
  \mathbf{F}_{HLL}=\frac{S_R\mathbf{F}_L-S_L\mathbf{F}_R+S_LS_R(\mathbf{U}_R-\mathbf{U}_L)}{S_R-S_L}\,.
  \end{equation}
  Robusto, disipativo. Es el recomendado para el código Python propio en la guía.
  \term{HLLD} agrega estados intermedios (mejor para contactos y ondas de Alfvén); fue
  el solver de la tarea de la onda de Alfvén.
\end{itemize}

\subsection{Reconstrucción y limitadores}

Para el Riemann hace falta el estado a cada lado de la cara. Con \term{reconstrucción}
\code{LINEAR} se interpola linealmente dentro de la celda; \code{PARABOLIC} (PPM) usa una
parábola (más precisa, la de la config \#07). El \term{limitador} (p.ej.\ \code{MC\_LIM})
evita que cerca de un choque la interpolación genere oscilaciones de Gibbs: recorta la
pendiente cuando el perfil tiene un salto. Es la consecuencia práctica del
\term{teorema de Godunov}: un esquema lineal de segundo orden no puede ser
monótono; hay que ser no lineal (limitar) para ser \term{TVD} (sin oscilaciones crecientes).

\subsection{Integración temporal y CFL}

El tiempo se avanza con \term{RK2}/\term{RK3} (Runge--Kutta de 2/3 etapas, método de
líneas) o el esquema de \term{HANCOCK} (predictor--corrector en un solo paso, el de la
config \#05). El paso de tiempo lo limita la condición \term{CFL}: la información no puede
recorrer más de una celda por paso, y como la onda más rápida es la magnetosónica rápida,
\begin{equation}
\Delta t = C_\mathrm{CFL}\,\min_\mathrm{celdas,dir}\Big(\frac{\Delta x}{|v_x|+c_{f,x}},
\ \frac{\Delta y}{|v_y|+c_{f,y}},\dots\Big),
\qquad C_\mathrm{CFL}\lesssim 0.3\text{--}0.8 .
\end{equation}
La corrida 3D usa CFL 0.25 (GLM añade la onda de limpieza a las velocidades, mejor
conservador). Físicamente: respetar el \term{dominio de dependencia} de las ecuaciones.

\subsection{Fronteras}

\term{Periódicas} en $x,z$: lo que sale por un lado entra por el otro (modo único bien
definido, sin bordes). \term{Reflectivas} en $y$: pared rígida, simetría de espejo; en la
práctica, en las celdas fantasma se copia el interior con \emph{signo invertido en las
componentes normales} de velocidad \emph{y} campo ($v_y\to-v_y$, $B_y\to-B_y$; lo demás
par). \term{Celdas fantasma} (\emph{ghost cells}): capas extra fuera del dominio que el
solver necesita para calcular flujos en el borde.

\section{Errores numéricos, Fourier y normas}
\label{sec:fourier}

\begin{itemize}[nosep]
  \item \term{Difusión numérica}: todo esquema disipa un poco; el efecto neto es como una
  viscosidad que suaviza gradientes y amontona el corte de la interfaz en
  $\sim$ pocas celdas. Por eso los modos con $\lambda\lesssim$ 4--8 celdas crecen menos de
  lo que dice la teoría, y por eso el $\sigma_\mathrm{eff}$ medido (0.229) quedó cerca del
  fundamental y no del modo de malla.
  \item \term{Nyquist}: una malla con espaciado $\Delta x$ no puede representar modos con
  $k>\pi/\Delta x$; en la corrida 3D, $\pi/\Delta x\approx151$. Es el ``corte de malla''
  que aparece en el mapa de teoría y la razón por la que el corte magnético duro
  ($k_c\approx157$) quedaba justo fuera del alcance.
  \item \term{FFT y espectro de potencia}: la transformada de Fourier discreta de un campo
  en la malla da la amplitud por modo; el \emph{espectro} es la potencia
  $P(\mathbf{k})=|\hat{q}(\mathbf{k})|^2$. En el notebook se usa sobre una lámina de
  densidad pegada a la interfaz, y se comparan \term{sectores} de ángulo alrededor de
  $k_x$ ($\parallel\B$) y $k_z$ ($\perp\B$): la razón $P_\perp/P_\parallel>1$ en los $k$
  altos es la firma espectral de la anisotropía magnética.
  \item \term{Normas} $L_1$ y $L_2$: error promedio (absoluto y cuadrático) entre dos
  campos, normalizadas por el área/volumen del dominio (y a veces por el promedio del
  campo, para hacerlas relativas). Se usan en el plan comparativo PLUTO vs.\ Python.
  \item \term{Convergencia}: repetir con mallas cada vez más finas; si el esquema es de
  orden $p$, el error cae como $(\Delta x)^p$; medir la pendiente en escala log--log da el
  ``orden observado''. Para la RT solo es limpio en fase lineal: después el sistema es
  caótico y dos corridas cualesquiera divergen punto a punto (por eso se comparan
  estadísticas, no campos).
\end{itemize}

\section{Los diagnósticos del notebook, explicados uno a uno}

\begin{itemize}[nosep]
  \item \term{Cortes 2D} ($z=0$ y $x=0$): proyección más simple; comparar los dos planos
  enseña la anisotropía (a lo largo de $B$ los dedos se ven gruesos y fundidos;
  transversal, finos y numerosos).
  \item \term{Promedios a lo largo de un eje}: $\langle\rho\rangle_x(y,z)$ elimina las
  estructuras alargadas en $x$ y deja la sección transversal de los dedos; simétrico en
  $z$. Es el truco estándar para ``ver'' hojas alargadas.
  \item \term{Espesor de mezcla $h_{95-5}$}: se buscan los cruces de
  $\rho=(\rho_1+\rho_2)/2$ en cada columna vertical y se toma la distancia entre los
  percentiles 95 y 5 de su distribución (percentiles y no max/min para que un solo dedo
  aventurero no infle la medida). Su ajuste log-lineal da $\sigma_\mathrm{eff}$.
  \item \term{Energías}: $E_\mathrm{int}=\int p/(\gamma-1)$, $E_\mathrm{cin}=\int\rho
  v^2/2$, $E_\mathrm{mag}=\int B^2/2$, $E_\mathrm{pot}=\int\rho\Phi$ con
  $\Phi=-\code{GRAV}\,y>0$ arriba (¡ojo con los signos: \code{GRAV} es negativo!). La
  conservación de la suma es el chequeo global del esquema conservativo + fuente.
  \item \term{Residuo de divergencia}: $\max|\nabla\cdot\B|\,\Delta x/B_0$ calculado con
  diferencias centradas del $B$ centrado en celdas; mide los monopolos de GLM.
\end{itemize}

%=====================================================================
\part*{Parte V --- Glosario alfabético rápido}
\addcontentsline{toc}{section}{Parte V --- Glosario alfabético}

\small
\begin{description}[style=nextline,leftmargin=1.2em,itemsep=1pt]
\item[AMR (\emph{adaptive mesh refinement})] refinamiento adaptativo de malla; PLUTO lo
tiene con Chombo; no se usa en nuestras corridas (malla estática).
\item[Atwood ($A$)] asimetría de densidades $(\rho_2-\rho_1)/(\rho_2+\rho_1)$; controla
tasas y velocidades asintóticas de burbujas/picos.
\item[Baroclínico] mecanismo que genera vorticidad cuando $\nabla\rho$ y $\nabla p$ no son
paralelos; es el motor de la RT.
\item[$B_c$, $\chi$] campo magnético crítico $\sqrt{(\rho_2-\rho_1)gL_x}$ y su fracción
usada ($B=\chi B_c/\sqrt{4\pi}$ en unidades de código). $\chi=1$ congela el modo
fundamental.
\item[Beta plasmática ($\beta$)] presión del gas relativa a la magnética, $2p/B^2$.
\item[CFL] condición que limita $\Delta t$ para que la información no cruce más de una
celda (Courant--Friedrichs--Lewy).
\item[Chevalier (1982)] solución auto-semejante eyecta (ley de potencias) + medio; la base
teórica de la fase en que el contacto del SNR es RT-inestable.
\item[Choque delantero / \emph{forward shock}] frente que comprime y calienta el medio
circundante; en él se drapea el campo ambiental.
\item[Chandrasekhar (1961)] libro clásico con la teoría lineal completa de la RT
magnetizada (relación de dispersión usada en el proyecto).
\item[CT (\emph{constrained transport})] actualización de $\B$ en malla escalonada con la
FEM; divergencia cero a precisión de máquina (config \#07 2D).
\item[CFL de GLM] GLM añade una onda rápida ($c_h$) al sistema: exige CFL más bajo
(0.25 en la 3D).
\item[Densidad conservativa vs.\ primitiva] $U=(\rho,\rho\vct,E,\B)$ vs.\
$V=(\rho,\vct,p,\B)$; el código evoluciona las primeras, los Riemann y la física se
leen en las segundas (conversión ida y vuelta constante).
\item[Discontinuidad de contacto] frontera entre eyecta chocada y medio chocado: presión
continua, densidad discontinua; ahí vive la RT del SNR.
\item[DivB, $\nabla\cdot\B$] divergencia magnética; debe ser cero (no hay monopolos);
su residuo numérico es un indicador de calidad del esquema.
\item[Drapeo] el campo ambiental se envuelve alrededor del choque/de los obstáculos,
ampliándose en los flancos.
\item[Eyecta] material expulsado por la supernova; en el mapeo es el fluido pesado.
\item[FEM (\emph{electromotive force})] $\mathbf{E}=-\vct\times\B$; en 2D solo hace falta
$E_z=(\vct\times\B)_z$, la que usa CT para actualizar $B_x,B_y$.
\item[FFT] transformada rápida de Fourier; base de los espectros del notebook.
\item[GLM (\emph{generalized Lagrange multiplier})] limpieza de divergencia de Dedner:
variable $\psi$ que propaga y amortigua el error (config \#09 3D).
\item[Godunov (método / teorema)] volúmenes finitos con solver de Riemann; el teorema
obliga a usar limitadores para segundo orden sin oscilaciones.
\item[HANCOCK] integrador predictor--corrector de un paso (config \#05).
\item[HLL / HLLD] solvers de Riemann de dos y cinco estados; HLLD es el preciso para MHD.
\item[Intercambio vs.\ undular] modos $\mathbf{k}\perp\B$ (sin freno magnético) vs.
$\mathbf{k}\parallel\B$ (estabilizados bajo $\lambda_c$).
\item[Kelvin--Helmholtz (KH)] inestabilidad de cizalla; secundaria en los flancos de
burbujas/picos; el campo la suprime en escalas cortas.
\item[$\lambda_c$, $\lambda_m$] longitud de corte magnético $\chi^2L_x$ y la del modo de
máximo crecimiento paralelo, $2\lambda_c$.
\item[Layzer / Goncharov] modelos de burbuja no lineal; $v_b=\sqrt{2Ag\lambda/[\pi(1+A)]}$
en 2D.
\item[Limitador MC] limitador de pendientes (monotonic-centered); evita oscilaciones.
\item[MHD ideal] fluido conductor perfecto; campo congelado; sin resistividad.
\item[Modo fundamental] el que cabe justo en el dominio, $\lambda=L_x$ ($k=2\pi/L_x$);
en 2D su tasa es $\sqrt{Agk}\sqrt{1-\chi^2}$.
\item[Monopolo] carga magnética neta espuria = error de $\nabla\cdot\B$; produce fuerzas
no físicas.
\item[Multimodo] perturbación con muchos $k$ (p.ej.\ la aleatoria 3D): capa de mezcla,
competencia de modos, estadística en vez de un solo pico.
\item[Nyquist ($k_N$)] número de onda máximo representable, $\pi/\Delta x$; el ``corte de
malla''.
\item[PLUTO] código de volúmenes finitos para astrofísica computacional (Mignone et al.\
2007, 2012); el usado en el proyecto.
\item[PPM / PARABOLIC] reconstrucción parabólica (piecewise parabolic method); más precisa
que LINEAR (config \#07).
\item[Relación de dispersión] ecuación que da $\sigma(k)$ (o $\omega(k)$); la de la RT
magnetizada es la eq.~\eqref{eq:sigma-mhd}.
\item[RK2 / RK3] Runge--Kutta de 2/3 etapas para el método de líneas.
\item[Roe] solver de Riemann por linealización de Roe; el de la corrida 3D.
\item[Sedov--Taylor] solución del punto explosivo en medio uniforme:
$R\propto(Et^2/\rho_0)^{1/5}$.
\item[SNR] supernova remnant = remanente de supernova.
\item[TVD] total variation diminishing: sin oscilaciones nuevas; propiedad que buscan los
limitadores.
\item[Unidades de código] variables adimensionales del solver; el setup fija
$L_x=\rho_1=c_{s,1}=1$.
\item[Vorticidad] $\boldsymbol{\omega}=\nabla\times\vct$; en la RT la genera el término
baroclínico en la interfaz.
\end{description}
\normalsize

%=====================================================================
\appendix
\section{Apéndice A --- Tabla de símbolos}

\begin{table}[H]
\centering\small
\caption{Todos los símbolos usados en el proyecto (¡atención a los choques de notación:
$A$ es el Atwood y también un área en las normas; $\eta$ es el cociente de densidades del
setup (\code{ETA}) y NO la resistividad, que no se usa en este problema!).}
\begin{tabular}{@{}lp{11.6cm}@{}}
\toprule
Símbolo & Significado \\
\midrule
$\rho$, $p$, $\vct$, $\B$ & densidad, presión, velocidad, campo magnético \\
$\gamma=5/3$ & razón de calor específico (gas monoatómico) \\
$E$ & energía total por volumen: interna $+$ cinética $+$ magnética \\
$e_\mathrm{int}=p/(\gamma-1)$ & energía interna por volumen \\
$\mathbf{g}$, $g$ & aceleración efectiva (constante, $\hat{\mathbf{y}}$, negativa) \\
$\Phi=-\code{GRAV}\,y$ & potencial gravitacional del término fuente \\
$e$, $\xi$ & (según contexto) energía interna específica / desplazamiento de la interfaz \\
$c_s$, $v_A$, $c_f$ & sonido, Alfvén, magnetosónica rápida \\
$\beta=2p/B^2$ & razón de presiones gas/campo \\
$A$ & número de Atwood $(\rho_2-\rho_1)/(\rho_2+\rho_1)$ \\
$\eta$ \code{(ETA)} & cociente de densidades $\rho_2/\rho_1$ del setup \\
$\chi$ \code{(CHI)} & campo en unidades del crítico: $B_0=\chi B_c/\sqrt{4\pi}$ \\
$B_c$ & campo crítico $\sqrt{(\rho_2-\rho_1)gL_x}$ \\
$k$, $\lambda$ & número de onda ($2\pi/\lambda$) y longitud de onda \\
$\sigma$ & tasa de crecimiento lineal ($e^{\sigma t}$) \\
$\lambda_c=\chi^2L_x$ & corte magnético (modos $\parallel\B$ menores quedan estables) \\
$\sigma_\mathrm{max}$ & tasa máxima de los modos paralelos, en $\lambda_m=2\lambda_c$ \\
$v_b$ & velocidad asintótica de burbuja (Layzer/Goncharov) \\
$h_{95-5}$ & espesor de la capa de mezcla (percentiles 95--5 de los cruces) \\
$U$, $V$ & variables conservativas y primitivas \\
$\mathbf{F}$ & flujo de las conservativas \\
$S_L$, $S_R$ & velocidades de las ondas extremas del solver de Riemann \\
$\mathbf{E}$, $E_z$ & campo eléctrico inducido (FEM) y su componente en 2D \\
$\psi$ & multiplicador de Lagrange de GLM (limpieza de div$B$) \\
$\Delta x$, $k_\mathrm{Nyq}$ & espaciado de malla y número de onda de Nyquist \\
$L_1$, $L_2$ & normas de error (promedio absoluto / cuadrático) \\
$P(\mathbf{k})$ & espectro de potencia $|\hat{q}|^2$ \\
\bottomrule
\end{tabular}
\end{table}

\section{Apéndice B --- Mapa: cada ecuación, dónde vive en el proyecto}

\begin{table}[H]
\centering\small
\caption{Dónde se usa cada resultado en los entregables.}
\begin{tabular}{@{}p{6.2cm}p{4.2cm}p{4.8cm}@{}}
\toprule
Ecuación / cantidad & Dónde se deriva o define & Dónde se usa \\
\midrule
Sistema MHD $+$ EOS & guía \S3; aquí Parte I & informe \S2; \code{pluto.ini} (la física del solver) \\
$c_s$, $v_A$, $c_f$, $\beta$ & guía \S3.2; aquí \S4 & CFL; diagnóstico del notebook \\
$\sigma^2=Agk$ (HD) & aquí \S6.3 & guía \S5.1; control 2D $\chi=0$ \\
$\sigma^2=Agk-2(\mathbf{k}\cdot\B)^2/(\rho_1+\rho_2)$ & aquí \S8 (\ref{eq:sigma-mhd}); guía \S5.2 & \code{teoria\_snr.py}; ajuste del notebook \\
$\lambda_c=\chi^2L_x$; $B_c$ & aquí \S8; guía \S5.3 & tabla de $\chi$ de la guía; mapa de teoría \\
$\sigma(\lambda)$, $\sigma_\mathrm{max}$, $\lambda_m$ & aquí \S8; \code{teoria\_snr.py} & interpretación de $h(t)$ y del espectro \\
$v_b$ (Layzer) & aquí \S7 & guía \S6; barrido de Atwood \\
$\xi(t)$ con semilla $v_0$ & aquí \S6.3 & cronología de la guía \S5.5 \\
Gravedad efectiva $|a|=\eta(1-\eta)R/t^2$ & aquí Parte III & informe \S2 (mapeo SNR) \\
CT / Faraday integral; GLM & aquí \S5; guía \S3.3 & config \#07 vs \#09; residuo divB \\
Flujo HLL; CFL & aquí Parte IV & guía \S8 (código Python propio) \\
Normas $L_1/L_2$ & aquí Parte IV; guía \S9 & plan comparativo del proyecto \\
\bottomrule
\end{tabular}
\end{table}

\end{document}
